%%%%%%%%%%%%%%%%%%%%%%%%%%%%%%%%%%%%%%%%%%%%%%%%%%%%%%%%%%%%%%%%%%%%%%%%%%%%%%%%
%   Copyright 2009, Oracle and/or its affiliates.
%   All rights reserved.
%
%
%   Use is subject to license terms.
%
%   This distribution may include materials developed by third parties.
%
%%%%%%%%%%%%%%%%%%%%%%%%%%%%%%%%%%%%%%%%%%%%%%%%%%%%%%%%%%%%%%%%%%%%%%%%%%%%%%%%

\section{\basiccore}\label{basic-calculus}

In this section, we define a basic core calculus for Fortress.
We call this calculus \emph{\basiccore}.
Following the precedent set by prior core calculi such as Featherweight
Generic Java~\cite{fgj},
we have abided by the restriction that all valid \basiccore\ programs are
valid Fortress programs.
\revision{revival-calculi}{This calculus predates instantiation exclusion
(\secref{trait-types}, \secref{trait-decls}).
Its only condition on multiple inheritance is that every visible method
have one owner, so an object that extends two different instantiations of
a parameterized trait declaring no method is a valid \basiccore\ program;
under the rule it is not a valid Fortress program, and the sentence above
no longer holds for every program.
The calculus, its rules and its soundness claim are as the team wrote and
proved them, and are not changed.}

\input{\calculi/basic/syntax}
\input{\calculi/basic/dynamic}
\input{\calculi/basic/static}
